\documentclass[11pt]{article}
\input{../calcnotes.sty}
\fancyhead[C]{1.3 Exponential Functions}

\begin{document}

\noindent An exponential function is defined as:

\vspace{1.5cm}

\noindent What are the two types of exponential functions? What do they look like?

\vspace{1.5cm}

\noindent Graph $y=2(3^x)-4$ and state the domain and range.

\emptygraph

\noindent Find the zero of $f(x)=5-2.5^x$ algebraically and graphically with a calculator.

\vspace{1.5cm}

\noindent Quick reminder of rules for exponents:

\begin{align*}
    x^5 \times x^2 = && \frac{x^5}{x^2} = && (x^5)^2 =\\
    (x^5 \times x^2)^3 = && \left(\frac{x^5}{x^2}\right)^3 = && \frac{1}{x^4} =
\end{align*}

\newpage

\noindent How do we model an exponential function? Model by hand and calculator:

\vspace{0.5cm}

\begin{tabular}{|c|c|}
    \hline
    \multicolumn{2}{|c|}{\textbf{U.S. Population}} \\ \hline
    \textbf{Year} & \textbf{Population (millions)} \\ \hline
    2002 & 287.9 \\ \hline
    2003 & 290.4 \\ \hline
    2004 & 293.2 \\ \hline
    2005 & 295.5 \\ \hline
\end{tabular}

\vspace{1.5cm}

\noindent What is the population today?

\vspace{1.5cm}

\noindent What other ways can we model exponential growth and decay? What is a half-life?

\vspace{1.5cm}

\noindent What is the yearly percent increase in US Population?

\vspace{1.5cm}

\noindent What is that number $e$? How can we approximate it with: $f(x)=\left(1+\frac{1}{x}\right)^x$

\end{document}
